\section{Reproducibility Audit}\label{app:repro}

Every numerical claim in this paper was re-derived from the stored per-fold prediction artifacts as a
final check before submission. We report the audit in full, including the places where it failed, because
a reader cannot otherwise distinguish numbers that have been verified from numbers that merely have not
been questioned.

\subsection{Protocol}
Each table cell was recomputed from the persisted subject-level prediction files
(\texttt{fold*/test\_eval/prism\_breath\_predictions.csv}) rather than copied from a logged summary, using
a single script (\texttt{scripts/regenerate\_supplementary\_tables.py}) that derives every column of a
table from one artifact and one label set. Where a reported number did not match, we searched
\emph{exhaustively}: all $66$ five-fold prediction sets stored under \texttt{outputs/} were scored, and
single runs, pairs and triples (for seed-averaged claims) were compared against the published values.
Because per-fold RR-MAE is independent of the labels, it serves as an artifact fingerprint that
identifies a run uniquely.

Two label conventions exist in the archive and are distinguished throughout. Nine neonates were stored
with a stale \texttt{age\_months}${=}6$ and were therefore scored against the $2$--$12$mo WHO cutoff
rather than the neonatal one. Correcting them moves the cohort from $86$ to $84$ tachypneic subjects and
the $<$2mo band from $29$ to $38$ subjects. All corrected-label results are reported here.

\subsection{What reproduces}
The audit reproduces the paper's central evidence exactly. Table~\ref{tab:audit} summarizes.

\begin{table}[!t]
\caption{Reproducibility audit. ``Exact'' means every cell matched the recomputed value to the printed
precision.}
\label{tab:audit}\centering\small
\setlength{\tabcolsep}{3.5pt}
\begin{tabular}{@{}llc@{}}\toprule
Table & Rows & Verdict\\\midrule
Objective isolation & all $9$ & Exact\\
Framing, salience rows & $2$ (three-seed each) & Exact\\
Framing, PRISM rate $\to$ threshold & $1$ & Exact\\
Component isolation, raw-flow FFT & $1$ & Exact\\
Ablations & $8$ of $9$ & Exact\\
Head-to-head, published extractors & $4$ & Exact\\
Per-fold supplement, published extractors & $4$ & Exact ($4$ d.p.)\\\midrule
Ablations, ``Full PRISM'' & $1$ & Artifact absent\\
Head-to-head, PRISM & $1$ & Artifact absent\\
Component isolation, fusion rows & $2$ & Artifact absent\\
Per-fold supplement, PRISM & $3$ & Artifact absent\\
Clinical operating points & all & Not uniquely identifiable\\
Radar and cross-dataset & all & Artifact absent\\
Independent replication & all & Artifact absent\\\bottomrule
\end{tabular}
\end{table}

The objective-isolation table---the frozen-encoder comparison that carries the paper's central
argument---matches its stored artifact row for row, including every RR-MAE and every DeLong $p$. The
three-seed salience results likewise reproduce exactly ($0.884{\pm}0.020$ and $0.899{\pm}0.035$), as do
all four published-extractor rows wherever they appear. The framing thesis is therefore verified
end-to-end.

\subsection{One missing artifact}
Every number that did \emph{not} reproduce traces to a single cause: the prediction file for the main
PRISM-Breath model run is no longer in the archive. No stored artifact reproduces its per-fold RR-MAE
within $0.10$ under either label convention, and the eight ablation rows that surround it in the main
ablation table all reproduce exactly---so the loss is specific to one run rather than systemic.

Restoring it is not simply a matter of retraining. No stored signal directory holds the learned-flow
representation together with depth for all $332$ subjects: the $332$-subject set carries the hand-crafted
flow, the learned-flow set covers $297$ subjects, and the $332$-subject learned set has its depth channel
zeroed. Reproducing the main run therefore requires re-extracting the learned flow for the $35$ absent
subjects before any training. An alternative that is reproducible today is to re-baseline the main model
on the $297$-subject learned-flow cohort, which matches the architecture described in the methods and the
cohort already used for the framing analysis.

The committed per-fold table additionally carries an internal inconsistency created by that loss: its
\texttt{Mean$\pm$SD} entry is the mean and \emph{population} standard deviation of the per-fold AUROCs
printed beside it, which correspond to the pre-correction labels, while its pooled cell reports the
post-correction value. The most likely explanation is that a label correction updated only the
label-dependent columns. We report these cells as they were computed, and flag them as not currently
regenerable. Re-executing the main configuration and regenerating the affected tables with the script
above is the single action that resolves all of them.

We note explicitly what this does \emph{not} affect. The comparison against published extractors is
unaffected in its baseline arm; the ablation conclusion (no single component is individually
load-bearing) rests on rows that all reproduce; and the objective and framing analyses do not use the
missing run.

\subsection{Numbers that are not independently verifiable}
Three further claims could not be checked and are marked as such rather than removed.

\textbf{Clinical operating points.} These are reported as a three-seed average. The averaging is not
uniquely identifiable from the archive: $56$ different triples of stored runs reproduce the full-cohort
values within tolerance, and no stored run is named for the seeds recorded in the hyperparameter table.
The values are plausible but should be regarded as unaudited until the seed artifacts are re-generated
and named.

\textbf{Radar and cross-dataset transfer.} No radar result artifact is retained; only the radar
processing script and the raw recordings remain, so the phantom and $50$-child results cannot be
recomputed here. One cross-dataset evaluation file survives and reports different values, but it is an
\emph{incomplete} run---its log terminates at $203$ of $252$ videos---so it neither confirms nor
supersedes the published rows. Both should be regenerated.

\textbf{Independent replication.} The external-cohort waveforms and their extraction code were not
retained, so those results cannot be recomputed here.

\subsection{Cohort and signal-set mismatches}
Two presentational issues surfaced that are independent of the missing artifact and that we correct in
the tables and captions.

First, the component-isolation table draws its raw-signal and encoder rows from the learned-flow
extraction on $297$ subjects while its fusion rows are on $332$; the reported gain is internally matched
but the column is not homogeneous. The raw-signal figure is also specific to that extraction: on the
alternative extraction over all $332$ subjects the same FFT read-out reaches AUROC $0.81$--$0.83$ rather
than $0.70$, so the low raw-signal baseline should not be read as a general property of optical flow.

Second, the framing table lists its salience rows ($297$ subjects) alongside a $332$-subject PRISM entry
in the same column. The statistics themselves are matched---the reported $p$ ranges are paired DeLong
tests against a $297$-subject regression-trained reference, and both the AUROCs and the $p$ ranges
reproduce exactly---and the caption states this. Only the visual juxtaposition of two cohorts in one
column invites misreading, which we resolve by labelling $n$ on every row.

\subsection{Corrections made during the audit}
The audit also surfaced implementation issues, all corrected in the released code.

\textbf{Calibration leakage.} Two analysis scripts selected a threshold or fitted a risk model on
validation predictions pooled across folds and then scored the pooled test set. Across this run the
validation splits draw from only $49$ distinct subjects, roughly ten of which fall in each test fold, so
the pooled procedure leaks. Both now calibrate per fold, and a guard raises rather than reporting a
result fitted on evaluated subjects.

\textbf{Conformal coverage.} The single pooled nonconformity width persisted in one artifact family
attains $0.762$ empirical coverage against a $0.80$ target, and $0.62$ for neonates. The age-conditional
rule described in the methods attains $0.81$ overall and $0.80$--$0.82$ in every band. Coverage is now
reported explicitly so the distinction is visible.

\textbf{Inactive audio branch.} The audio branch receives zero fusion weight for all $332$ subjects and
its counterfactual removal changes the predicted rate by exactly zero for every subject. The system is
described as two-modality throughout.

\textbf{Near-uniform fusion weights.} The learned reliability weights are $0.498{\pm}0.007$ (depth) and
$0.502{\pm}0.007$ (flow), so the fusion behaves as an equal-weight average. This is now stated where the
weighting is defined.

\subsection{Release}
The audit script, the table regeneration script, and the per-subject outputs behind every table in this
appendix are released with the code. A reader who re-runs the regeneration script against the archive
will reproduce the verdicts in Table~\ref{tab:audit}, including the failures.
